%%%%%%%%%%%%%%%%%%%%%%%%%%%
\section{Conclusioni}
%%%%%%%%%%%%%%%%%%%%%%%%%%%%%%%%%%%%%%%%%%%%%%%%%%%%%%%%%%%%%%%%%%%%%%%%%%%%%%%%%%%%%%%%%%%%%%%%%%%%%%%%%%
\begin{frame}
\vspace{0.5cm}
\frametitle{Conclusioni}
\begin{block}{Risultati Ottenuti}
\begin{itemize}
\item L'evoluzione della composizione del combustibile è in linea con altri studi
\item L'approccio multifisico consente di effettuare una simulazione complessiva della fisica del reattore
\item L'accoppiamento neutronico-termoidraulico ottenuto appare promettente e valido

% si è presentato valido, come mostrato dai risultati
%       soddisfacenti

% apre a prospettive future per l'indagine del comportamento multifisico
% di sistemi nucleari

%\item La soluzione è ottenuta dopo numerose iterazione attraverso il codice di neutronica e termoidraulica con correzioni di temperatura e densità
%\item I risultati ottenuti aprono a prosfettivi di estendere un simile approccio ad altri modelli
%\item L'intera procedura è in grado di garantire flessibilità all'utente
\end{itemize}
\end{block}
\begin{block}{Sviluppi Futuri}
\begin{itemize}
\item Aggiornamento dei dati utilizzati nella modellazione con i progetti più recenti inerenti al progetto ALFRED
\item Approfondimento sui sistemi di sicurezza (barre di controllo e di sicurezza) del reattore nel codice di neutronica
\item Raffinazione del modello termoidraulico attraverso l'introduzione di modelli più sofisticati 
      e l'implementazione di codici per la turbolenza dei metalli liquidi
% \item Estendere un simile approccio multifisico ad altri reattori
% \item $k_{eff}$ troppo elevato in riferimento ad altri studi del settore
% \item Possibilità di interpolare solo su due variabili e non su tre 
% \item Modello di termoidraulica troppo semplificativo per un'analisi approffondita
\end{itemize}
\end{block}

\end{frame}

% \section{Sviluppi Futuri}
%%%%%%%%%%%%%%%%%%%%%%%%%%%%%%%%%%%%%%%%%%%%%%%%%%%%%%%%%%%%%%%%%%%%%%%%%%%%%%%%%%%%%%%%%%%%%